\documentclass[]{article}
\usepackage{graphicx}
\usepackage{amsmath,amsfonts}
\usepackage{kerkis}
\usepackage[english,greek]{babel}
\usepackage[dvipsnames]{xcolor}
\usepackage{xcolor}
\usepackage{multirow,bigdelim}
\usepackage{caption}

%opening
\title{}
\author{}

\begin{document}

\maketitle

\begin{abstract}

\end{abstract}
\section{Απόδειξη 1}
\textbf{Άσκηση:} Το δυναμοσύνολο του $\mathbb{N}$ είναι μη-αριθμήσιμο.
Αποδείξτε.\\
\textbf{Λύση:} Έστω $\mathcal{P}(\mathbb{N})$ αριθμήσιμα άπειρο. Δοθέντος ενός συνόλου $A$, έχουμε:
\begin{center}
	Έστω $A = \{1, 2\}$ και\\
	$\mathcal{P}(A) = \{\emptyset, \{1\}, \{2\}, \{1, 2\}\}$\\
\end{center}
Τα στοιχεία του δυναμοσυνόλου του Α μπορούν να απαριθμηθούν. Τοιουτοντρόπως, μπορούμε να απαριθμησουμε \& τα στοιχεία του $\mathcal{P}(\mathbb{N})$. \\
Έστω δυαδική σχέση $R \subseteq (A \times A)$. Σύμφωνα με την αρχή της διαγωνοποιήσης, το διαγώνιο σύνολο $D$ είναι το σύνολο που ορίζεται ως
$D = \{a : a \in A \land (a,a) \notin R\}$. Σύμφωνα με την ίδια αρχή, τα σύνολα $R_a$ ορίζονται ως
$R_a = \{b : b \in A \land (a,b) \in R\} \rightarrow$, όπου το $R_a$ δηλαδή περιλαμβάνει κάθε $b$ για το οποίο ισχύει ότι το διατεταγμένο ζεύγος $(a,b)$ ανήκει στην δυαδική σχέση $R$. Π.χ.
$A = \{0, 1, 2\}$, $R = \{(0, 1), (1, 0)\}$
\begin{center}
	\begin{tabular}{c|cc}
		& 0 & 1 \\ \hline
		0 &   & + \\ \hline
		1 & + &   
	\end{tabular}
\end{center}
Για το παραπάνω παράδειγμα έχουμε $R_0 = \{1\} \rightarrow \exists (0,1)$ και $R_1 = \{0\}$ γιατί $\exists (1,0)$. Δεδομένων των προαναφερθέντων, δύναται να ορίσουμε μια δυαδική σχέση ως εξής:
\begin{equation}
R = \{(i,j) : j \in R_i\} \iff R_i = \{j : j \in \mathbb{N} \land (i,j) \in R\}
\end{equation}

Από τη στιγμή που ένα σύνολο δεν αποτελεί μια αυστηρή διάταξη θα μπορούσαμε να έχουμε μια αναπαράσταση της δυαδικής μας σχέσης ως εξής:

\begin{center}
	\captionof{table}{Φανερό Τμήμα Δυαδικής Σχέσης $R$}
	\begin{tabular}{rl|c|c|c|c|c|c|}
		\multicolumn{2}{c}{} & \multicolumn{6}{c}{\makebox[0pt]{\footnotesize Αύξων Αρ.}} \\[-1ex]
		\multicolumn{2}{c}{} & \multicolumn{6}{c}{\downbracefill} \\
		\cline{3-8}
		& & 0 & 1 & 2 & 3 & 4 & 5 \\ \cline{2-8}
		\ldelim\{{6}{*}[{\footnotesize Νούμερο }] & \multicolumn{1}{|l|}{0} &   & + &   &   &   &   \\ \cline{2-8}
		& \multicolumn{1}{|l|}{1} & + &   & + &   &   &   \\ \cline{2-8}
		& \multicolumn{1}{|l|}{2} &   & + &   &   & + &   \\ \cline{2-8}
		& \multicolumn{1}{|l|}{3} &   & + &   &   & + &   \\ \cline{2-8}
		& \multicolumn{1}{|l|}{4} &   &   &   &   & + &   \\ \cline{2-8}
		& \multicolumn{1}{|l|}{5} &   &   &   &   &   &   \\ \cline{2-8}
	\end{tabular}
\end{center}
 Έχω δυαδική σχέση $R = \{(i,j)$ αν $j \in R_i\}$. Δηλαδή αμα έχω $(1, 0)$, $R_1 = \{0, 2\}$, $(1, 2)$. Ύστερα, σημειώνουμε τα υπάρχοντα στοιχεία \& αυθαίρετα τα υπόλοιπα.


Ο πίνακας αυτός δομείται ως $N_0, N_1, N_2 \dots N_x \in \mathcal{P}(\mathbb{N})$, καθώς αν είναι πραγματικά αριθμήσιμο (ή αριθμήσιμα άπειρο, ορθότερα) θα πρέπει να υπάρχει τρόπος να δεικτοδοτήσουμε τον κάθε πιθανό συνδυασμό-υποσύνολο που μπορεί να βρεθεί, έστω και με τυχαία διαδικασία.
Γενικά, δεν χρειάζεται να είναι ο δείκτης ίσο με μέγεθος υποσυνόλου. Ειδικότερα, ο δείκτης, καθώς υποθέτουμε ότι είναι αριθμήσιμο, ανήκει στο $\mathbb{N}$ και αποτελεί ένα index που αντιστοιχεί σε κάθε μέλος του $\mathcal{P}(\mathbb{N})$ δηλ.
\begin{equation}
	R = \{(i, j) : j \in N_i \land N_i \in \mathcal{P}(\mathbb{N}) \land i \in \mathbb{N}^0\}
\end{equation}


Έστω αυθαίρετα επιλεγμένα $N_i \in \mathcal{P}(\mathbb{N})$, $i \in \{0..4\}$, με τις υπόλοιπες τιμές τους \& μέλη του $\mathcal{P}(\mathbb{N})$ να μην μας αφορούν.
Ας λάβουμε τώρα ως $D$ το προαναφερθέν $D$, δηλ. όπου στην κύρια διαγώνιο σημειώνεται στοιχείο θα είναι κενή η αντίστοιχη θέση \& με (+) αν είναι αρχικά κενή.
Αν $R_D = \{b : b \in \mathbb{N} \land (b,b) \in R\}$ τότε $D = \overline{R_D} \Rightarrow$
$ D = \{0, 1, 2, 3 \dots\}$. Ας το σχεδιάσουμε , έστω, σε μορφή στήλης.
\begin{center}
	\begin{tabular}{c}
		$D$ \\ \hline
		+ \\ \hline
		+ \\ \hline
		+ \\ \hline
		\textcolor{red}{+} \\ \hline
		\\ \hline
		$\vdots$
	\end{tabular}
\end{center}
Οφείλουμε τώρα να σκεφτούμε ότι το $D$ συνίσταται από $n \in \mathbb{N}$, άρα θα πρέπει να βρεθεί μέσα στην απαρίθμηση των μελών του $\mathcal{P}(\mathbb{N})$ $\{R_0, R_1 \dots\}$ ως υποσύνολο όλου του συνόλου των φυσικών αριθμών π.χ. $\{1, 2, 3, 4\} \subseteq \mathbb{N}$ άρα $\{1, 2, 3, 4\} \in \mathcal{P}(\mathbb{N})$ βάσει λογικής. Ας δούμε τι ισχύει στην πράξη. \\


	 Ας δούμε αν είναι ίδιο με το $R_0$. Είναι? Όχι! Γιατί? $0 \notin R_0, 0 \in D$. 
	 Μήπως αν δείξουμε ότι το $D$ ταυτίζεται με το $R_1$? Μπα! Εδώ θα διαφέρει καθώς $1 \notin R_1, 1 \in D$. Για το $R_2$? Τσου, έχουμε $2 \notin R_2, 2 \in D$ κ.ο.κ.
	 Γενικά, βλέπουμε ότι διαφέρει από κάθε $R_i$: πράγμα που σημαίνει ότι δεν είναι εντός της $R$.
	 Όμως, εμείς είπαμε ότι "αριθμήσιμα άπειρο, άρα μπορώ να συνδέσω \& βρω κάθε πιθανό στοιχείο του συνόλου σαν φύση με δείκτη"
	 Από τη στιγμή που κάθε [αριθμήσιμο] βρίσκεται εκτός της λίστας μας, τότε δεν είναι αριθμήσιμο άρα καταλήγουμε σε άτοπο. Δηλαδή, το $\mathcal{P}(\mathbb{N})$ \textbf{δεν είναι αριθμήσιμα άπειρο}!

\section{Απόδειξη 2}
\textbf{Άσκηση:} Αποδείξατε ότι το σύνολο όλων των πραγματικών αριθμών στο διάστημα $[0,1]$ \textbf{δεν} είναι μετρήσιμο.
\\~\\
\textbf{Λύση:}
Όπως γνωρίζουμε, ένας πραγματικός αριθμός μπορεί να έχει άπειρα ψηφία τα οποία πολλές φορές αγνοούμε καθώς είναι μηδενικά μέχρι το τέλος. Ωστόσο, υπάρχουν και περιπτώσεις όπου δύναται να τελειώνει με κάποιο ψηφίο.
Εδώ, ας υποθέσουμε ότι ΕΊΝΑΙ μετρήσιμο το σύνολο των πραγματικών αριθμών μεταξύ 0 και 1, \textbf{συμπεριλαμβανομένων αυτών}. Αυτό σημαίνει ότι μπορούμε να τους κατατάξουμε σε μια λίστα με κάποιον φυσικό αριθμό χωρίς κατα ανάγκη να βάζουμε τον μικρότερο πρώτο λόγω του ότι α) χειριζόμαστε σύνολα και β) η λίστα , ή καλύτερα, απαρρίθμηση των πραγματικών μεταξύ 0 και 1 , συμπεριλαμβανομένων αυτών, δομείται από εμάς.

Για ένα πραγματικό αριθμό μεταξύ του 0 και του ενός έχουμε τα $ \{0\} \in \mathbb{Z} \subset \mathbb{R} $ και $ \{1\} \in \mathbb{Z} \subset \mathbb{R} $. Ωστόσο έχουμε και τους ενδιάμεσους αριθμούς άρα $(0,1)$. Επομένως $[0,1] = \{0\} \cup (0,1) \cup \{1\} \rightarrow [0,1] = [0,1) \cup \{1\} $. Αν το [0,1] είναι αριθμήσιμο, τότε θα πρέπει και το $[0,1)$ να είναι αριθμήσιμο καθώς συντελεί ένα αριθμήσιμο σύνολο (μαζί με το σύνολο $\{1\}$). \\
Θυμηθείτε ότι υπάρχουν άπειρα δεκαδικά ψηφία , έστω εν δυνάμει. Ωστόσο, αν θέλουμε να υποθέσουμε ότι το [0,1] είναι αριθμήσιμα άπειρο, τότε θα πρέπει να είναι και το (0,1). Αυτό σημαίνει, με τη σειρά του, ότι θα πρέπει να υπάρχει η ικανότητα να κατατάξουμε τα στοιχεία του συνόλου αυτού με τη βοήθεια κάποιου δείκτη. 

Ας θεωρήσουμε λοιπόν όλα μαζί πως το σύνολο [0,1) δομείται ως
\begin{equation*}
	 P = \{p_0,p_1,p_2,...\}
\end{equation*}, όπου $p$ ένας δεκαδικός αριθμός με μηδενικό ακέραιο μέρος. Σημειώνουμε επίσης τη μεταβλητή $k_x^y: P_x \in P$ , η οποία λαμβάνει τιμή εκ της κάτωθι συνάρτησης (μετ''εμπνεύσεως από τις διαφάνειες του κ.Σπηλιόπουλου): 
\begin{equation}
	k_x^y =
	\begin{cases}
		1 & \text{, αν το $y+1$-οστό ψηφίο του $p_x$ \textbf{δεν} είναι ίσο με } 0 \\
		0 & \text{, αν το $y+1$-οστό ψηφίο του $p_x$ είναι ίσο με }  0
	\end{cases}
\end{equation}
  , με το $y+1$-οστό ψηφίο του $p_x$ προς έλεγχο να απορέει από την επόμενη σχέση.
  \begin{equation}
  	d_x^y = \left\lfloor\frac{p_x}{10^{-(y+1)}}\right\rfloor\mod{10}
  \end{equation} Σε αυτό το σύνολο λοιπόν θα αφορά η δυαδική σχέση $R=\{(a,b): b\in R_a\}$, όπου για το $R_a$ έχουμε $R_a=\{b: b \in \mathbb{N} \wedge k_a^b = 1, p_a \in P \}$. Δηλαδή το δεκαδικό ψηφίο του $a$-οστού αριθμού στην $(b+1)$-οστή θέση να μην είναι 0. Ένας τρόπος να απεικονίσουμε αυτή τη σχέση , έστω με τυχαίες τιμές να φαίνονται πρώτα, είναι ο εξής:  
%  \begin{eqnarray*}
%  	\left.
%  	\begin{array}{ll}
%  		\mbox{Aardvark} & \mbox{Armenia}\\ 
%  		\mbox{Elephant} & \mbox{Ethiopia} 
%  	\end{array} 
%  	\right\}  && \mbox{Things beginning with vowels} \\ 
%  	\mbox{Things beginning with consonants:} && \left\{
%  	\begin{array}{ll}
%  		\mbox{Platypus} & \mbox{Portugal}\\ 
%  		\mbox{Zebra} & \mbox{Zimbabwe} 
%  	\end{array}
%  	\right.
%  \end{eqnarray*}

\begin{center}
\captionof{table}{Παραδειγματικά Φανερό Απόσπασμα της $R$}

\begin{tabular}{rl|l|l|l|l|l}
	\multicolumn{2}{c}{} & \multicolumn{4}{c}{\makebox[0pt]{\footnotesize Θέση Ψηφίου}} \\[-1ex]
	\multicolumn{2}{c}{} & \multicolumn{4}{c}{\downbracefill} \\
	\cline{3-6}
	& & 0 & 1 & 2 & 3 \\ \cline{2-6}
	\ldelim\{{4}{*}[{\footnotesize Αριθμός $p_a$ }] & \multicolumn{1}{|l|}{0} & + &   & + & + & ... \\ \cline{2-6}
	& \multicolumn{1}{|l|}{1} &   &   & + & & ...  \\ \cline{2-6}
	& \multicolumn{1}{|l|}{2} &   &   &   & & ... \\ \cline{2-6}
	& \multicolumn{1}{|l|}{3} &   &   & + & + & ... \\ \cline{2-6}
	\multicolumn{6}{r}{}  ...............................    \\
	\multicolumn{6}{r}{}    ...............................
\end{tabular}

\end{center}
Έστω $D : \{ b : b \in \mathbb{N} \wedge (b,b) \notin R \}$ βασει της αρχής της διαγωνιοποιήσης.
Εδώ μιλώντας βάσει του φανερού μέρους της $R$, αν για $A = \{0,1,2,3\}$ έχουμε σημειωμένα τα $\{0,3\}$ στην διαγώνιο την κύρια δηλαδή για $\Delta = \{ c : c \in \mathbb{N} \wedge (c,c) \in R \}$ έχουμε $\{0,3\} \in \Delta$ και $\{1,2\} \notin \Delta$ τότε $\{0,3\} \notin D$ και $\{1,2\} \in D$.

Θα δείτε λοιπόν οτι αν πάρουμε το $D$ (διαγώνιο σύνολο) σαν έναν κωδικοποιημένο αριθμό\footnote{
Ο αριθμός αυτός παράγεται με τον απλό τύπο $D_{coded} = \sum_{i \in \mathbb{N}} D_i * 10^{-i-1}$ 
} , θα πρέπει να ανήκει στο $[0,1]$ καθώς, σαν ενα αριθμήσιμο σύνολο, οφείλει αυτο να μπορεί να "δείξει" το $D$ με κάποιο δείκτη (δηλ. να ανήκει στην απαρρίθμηση του).

Ωστόσο, αν προβούμε στην σύγκριση του $D$ σαν αριθμό - γραμμή με κάθε $R_a$ , βλέπουμε το εξής:
\begin{itemize}
	\item Συγκρίνοντας με τον $1^\text{ο}$ κωδ/μενο αριθμό, ο $D$ \underline{δεν} περιλαμβάνει μη-μηδενικό ψηφίο στην $(0+1)$η θέση, ενω ο $p_0$ ναι! Άρα διαφέρουν!
	\item Συγκρίνοντας με τον $2^\text{ο}$, θα δω οτι ο $p_1$ εχει $k_1^1 = 0$, ενω ο $D$ ναι, άρα διαφέρουν.
\end{itemize}
Αν συνεχίσουμε, θα δούμε ότι το $D$ διαφέρει απο κάθε στοιχείο της απαρρίθμησης. \\
Εμείς ομως, υποθέσαμε οτι κάθε μέλος του $[0,1)$ θα μπορούσαμε να το βρούμε σα να διατρέχουμε μια λίστα. Αντ'αυτού, αφού $D \neq p_0, D \neq p_1$ κ.ο.κ.  φαίνεται ότι δεν μπορεί να βρεθεί (καθώς δεν ταυτίζεται με τα υπάρχοντα στοιχεία).

Αυτό βρίσκεται σε πλήρη αντίθεση με την υπόθεση μας οτι το $[0,1)$ θα είναι αριθμήσιμο αρα καταλήγουμε σε \textbf{άτοπο}, δείχνοντας έτσι ότι το $[0,1)$ είναι μη-μετρήσιμο. Από τη στιγμή που δεν είναι αριθμήσιμο , τοτε και το $[0,1]$ , ως ένωση του συνόλου αυτού με το $\{1\}$ , θα είναι και αυτό \textbf{μη αριθμήσιμο!} 
\\~\\
\textit{Ερώτηση: ΑΝ το κάναμε σε δυαδικό σύστημα, πόσο πιο εύκολο θα ήταν?}
\pagebreak
\section{Απόδειξη 3 (ΣΗΜΑΝΤΙΚΉ)}
\textbf{Άσκηση}:  Το σύνολο των γλωσσών ενός αλφαβήτου $\Sigma$ είναι μη αριθμήσιμο. Να αποδειχθεί η προηγούμενη πρόταση.
\textbf{Λύση}: 

Έστω $w$ συμβολιζόμενες οι συμβολοσειρές όπως τις μάθαμε \& $L$ οι γλώσσες. Ως $L^A$ ορίζουμε το σύνολο όλων των γλωσσών επί του αλφαβήτου $\Sigma$. Υποθέτοντας ότι το $L^A$ είναι αριθμήσιμο, αυτό σημαίνει ότι μπορούμε να απαριθμήσουμε τα στοιχεία του ως
$$L^A = \{L_0, L_1, L_2 \dots\}$$ άρα:
\begin{equation*}
	L^A = \{L_i : i \in \mathbb{N} \land L_i \subseteq \Sigma^*\}
\end{equation*}
Με τη σειρά του, αν θυμηθούμε ότι το $\Sigma^*$ είναι το σύνολο των συμβολοσειρών επί ενός αλφαβήτου $\Sigma$ έχουμε
\begin{equation*}
\Sigma^* = \{w_j : j \in \mathbb{N} \text{ και η } w_j \text{ συντίθεται από σύμβολα του } \Sigma\}
\end{equation*}
Ας χτίσουμε τώρα μια δυαδική σχέση $R \subseteq \mathbb{N} \times \mathbb{N}$ τέτοια ώστε να υπάρχει ένας δείκτης $i$ άρα να δεικτοδοτείται σημασιολογικά σε ποιά γλώσσα βρισκόμαστε , και αντίστοιχα ένας τέτοιος δείκτης $j$ για το για ποια συμβολοσειρά μιλάμε. Έχουμε δηλαδή
$$R = \{(a,b) : a,b \in \mathbb{N} \land L_a \in L^A \land b \in R_a\} ~ \text{με} $$
$$R_a = \{b : b \in \mathbb{N} \land w_b \in L_a\}$$
Παρατηρήστε το παραδειγματικά ορατό τμήμα της σχέσης που έπεται της πρότασης.
\begin{center}
		\captionof{table}{Ορατό Τμήμα Δυαδικής Σχέσης $R$}
	\begin{tabular}{rl|c|c|c|}
		\multicolumn{2}{c}{} & \multicolumn{3}{c}{\makebox[0pt]{\footnotesize Συμβολοσειρά (Δείκτης συμ/σειράς)}} \\[-1ex]
		\multicolumn{2}{c}{} & \multicolumn{3}{c}{\downbracefill} \\
		\cline{3-5}
		& & 0 & 1 & 2 \\ \cline{2-5}
		\ldelim\{{3}{*}[{\footnotesize Γλώσσα (Δείκτης γλ.) }] & \multicolumn{1}{|l|}{0} &   & + &   \\ \cline{2-5}
		& \multicolumn{1}{|l|}{1} &   & + &   \\ \cline{2-5}
		& \multicolumn{1}{|l|}{2} &   &   &   \\ \cline{2-5}
	\end{tabular}
	\quad $R_0 = \{1\}$ άρα $w_1 \in L_0$
\end{center}

Ορίζουμε ως $D$  ως $D = \mathbb{N} - R_\Delta$ όπου $R_\Delta = \{a : a \in \mathbb{N} \land (a,a) \in R\}$. Για το ορατό κομμάτι του παραδείγματος, αν $1 \in R_\Delta$ και $0, 2 \notin R_\Delta$ τότε $1 \notin D$ και $0, 2 \in D$.

Έπειτα, ας θυμηθούμε το εξής: η $L^A$ σαν απαρίθμηση όλων των γλωσσών $L \subseteq \Sigma^*$, θα πρέπει να έχει φροντίσει να λιστάρει και την γλώσσα που σημειώνεται από το διαγώνιο σύνολο καθώς για κάθε $x \in D : w_x \in \Sigma^*$. Ας πάρουμε το παραδειγματικό κομμάτι. \pagebreak

\begin{center}
	\captionof{table}{Ορατό Τμήμα Διαγώνιου Συνόλου $D$}
	\begin{tabular}{|c|}
		\multicolumn{1}{c}{$D$} \\ \hline
		+ \\ \hline
		\\ \hline
		+ \\ \hline
		\multicolumn{1}{c}{$\vdots$}
	\end{tabular}
\end{center}

Αν όπου (+) σημειώνεται η ύπαρξη λέξεως, τότε $w_0, w_2 \in L_D$.
Ας δούμε λοιπόν αν ταυτίζεται με κάποια προϋπάρχουσα γλώσσα:
\begin{itemize}
	\item Η $L_D$ διαφέρει από την $L_0$ καθώς $w_0 \notin L_0$ αλλά $w_0 \in L_D$
	\item Η $L_D$ διαφέρει από την γλώσσα $L_1$ καθώς $w_1 \in L_1$ αλλά $w_1 \notin L_D$.
	\item Η $L_D$ διαφέρει από την γλώσσα $L_2$ ($w_2 \notin L_2$ αλλά $w_2 \in L_D$).
\end{itemize}

Βλέπουμε ότι, με τη βοήθεια της αρχής της διαγωνιοποίησης, ότι καμία γλώσσα της $L^A$ δεν ταυτίζεται με την $L_D$ άρα δεν έχει συμπεριληφθεί στην απαρίθμηση. Ωστόσο, σαν αριθμήσιμα άπειρο σύνολο, θα έπρεπε το σύνολο όλων των γλωσσών να επιτρέπει την απαρίθμηση όλων των στοιχείων του ως αμφιμονοσήμαντη συνάρτηση με το $\mathbb{N}$ (δηλ. λαμβάνοντας υπόψη \textbf{και} την $L_D$). Έτσι, οδηγούμαστε σε \textbf{άτοπο}, καθιστώντας την αρχική υπόθεση μας περί αριθμησιμότητας του $L^A$ λάθος.

Συμπερασματικά, αποδείχθηκε ότι το σύνολο όλων των γλωσσών επί ενός αλφαβήτου $\Sigma$ είναι μη-αριθμήσιμο!
\end{document}
